% ---------------------------------------------------------------------------
% Compact print layout — PDF output ONLY (loaded from `format: pdf` in
% _quarto.yml, so the HTML book is untouched).
%
% The handout is printed and handed to students, so every saved page is real
% paper. Everything below buys pages back from *whitespace*, never from the
% text itself: the body font, the line length and the leading stay in the
% comfortable range, and only the padding around headings, lists, code blocks,
% callouts and figures gets trimmed.
%
% This file is loaded through `include-in-header`, which lands in the preamble
% AFTER pandoc's syntax-highlighting setup (fvextra) and after Quarto's own
% tcolorbox definitions for code blocks and callouts — so `\fvset` and
% `\tcbset` below can safely re-tune both.
% ---------------------------------------------------------------------------

% --- Justification ---------------------------------------------------------
% A compact page has long lines, and long lines are where justification goes
% wrong: TeX stretches the spaces of one line and squeezes the next.
% microtype lets punctuation hang a hair into the margin and tunes the space
% around it, which evens the grey of the text out again. It also happens to
% fit a few more words per line.
\usepackage{microtype}

% --- Paragraph spacing -----------------------------------------------------
% Pandoc's KOMA branch sets `parskip=half`: no first-line indent, half a blank
% line between paragraphs. The `-` variant drops KOMA's extra requirement that
% the last line of a paragraph must leave a visible gap, which alone removes a
% forced line here and there. The explicit length then goes from half a line
% to about a third — still an obvious break between paragraphs, but roughly
% two lines per page cheaper.
\KOMAoptions{parskip=half-}
\AtBeginDocument{\setlength{\parskip}{.35\baselineskip plus .1\baselineskip}}

% --- Heading spacing -------------------------------------------------------
% KOMA is generous above and below every heading. The values below keep a
% clear visual break (a heading must still "belong" to the text under it) at
% about two thirds of the original space.
%
% `beforeskip` is negative on purpose for the section levels: the minus sign
% is KOMA's flag for "do not indent the first paragraph after this heading",
% not a negative distance.
%
% A chapter always starts on a fresh page, so the 2.3 blank lines KOMA puts
% *above* the chapter title are pure loss — dropping them saves a few lines in
% each of the ~45 chapters.
\RedeclareSectionCommand[beforeskip=0pt,afterskip=1\baselineskip]{chapter}
\RedeclareSectionCommand[
  beforeskip=-2.2ex plus -.5ex minus -.2ex,
  afterskip=.9ex plus .1ex]{section}
\RedeclareSectionCommand[
  beforeskip=-1.9ex plus -.5ex minus -.2ex,
  afterskip=.7ex plus .1ex]{subsection}
\RedeclareSectionCommand[
  beforeskip=-1.7ex plus -.5ex minus -.2ex,
  afterskip=.6ex plus .1ex]{subsubsection}

% --- Lists -----------------------------------------------------------------
% The chapters use a lot of numbered exercise steps and bullet lists. Pandoc's
% \tightlist already removes the space *between* items of a tight list, but
% not the space above and below the list itself, and it cannot: that space is
% consumed before \tightlist runs. enumitem sets it for every list at once.
\usepackage{enumitem}
\setlist{topsep=2pt,partopsep=0pt,itemsep=2pt,parsep=3pt}

% --- Code blocks and callouts ----------------------------------------------
% Both are tcolorbox environments (Quarto builds them that way), so one
% \tcbset tightens all of them. Quarto sets `left` itself for callouts, so
% that value is left alone here and only the vertical padding shrinks.
% \fvset controls the verbatim text inside a code block: one step smaller than
% the body text, and without the extra leading the body text gets — code is
% read line by line, so it does not need the same open spacing.
\AtBeginDocument{%
  \tcbset{%
    boxsep=0mm,
    top=1mm, bottom=1mm,
    before skip=1.2ex, after skip=1.2ex,
  }%
  \fvset{fontsize=\small,baselinestretch=1}%
}

% --- Figures ---------------------------------------------------------------
% LaTeX defaults are conservative: it refuses to put a figure on a page that
% is more than 70% text, and happily banishes a big figure to a page of its
% own. The diagrams in this book are small enough to sit inside the text, so
% allow much fuller pages and only fall back to a float page when a figure
% really fills one.
\renewcommand{\topfraction}{0.9}
\renewcommand{\bottomfraction}{0.9}
\renewcommand{\textfraction}{0.06}
\renewcommand{\floatpagefraction}{0.8}
\setlength{\intextsep}{6pt plus 2pt minus 2pt}
\setlength{\textfloatsep}{8pt plus 2pt minus 2pt}
\setlength{\abovecaptionskip}{4pt}
\setlength{\belowcaptionskip}{0pt}
